\documentclass[../../../include/open-logic-section]{subfiles}

\begin{document}

\olfileid[fr]{sth}{z}{arbintersections}
\olsection{Annexe : clôture, compréhension et intersection}

Dans la \olref[sth][z][milestone]{sec}, nous vous avons proposé de
revenir sur le travail naïf de la \olref[sfr][][]{part} pour vérifier
qu'on peut le mener dans~$\Zminus$. Si vous avez suivi ce conseil,
\emph{un} point a pu vous arrêter : l'usage de
l'\emph{intersection} dans le traitement des \emph{clôtures} par
Dedekind.

La \olref[sfr][infinite][dedekind]{Closure} fixe un ensemble ambiant.
La clôture est l'intersection de ses parties stables par $f$ qui
contiennent $o$. Voici la formulation dans laquelle l'ensemble ambiant
reste implicite :
\begin{align*}
  \closureofunder{f}{o} & = \bigcap\Setabs{X}{o \in X \text{ et $X$ est
stable par $f$}}.
\intertext{La forme générale est celle d'une définition comme :}
  C & = \bigcap\Setabs{X}{\phi(X)}.
\end{align*}
Mais cette forme devrait éveiller nos soupçons : puisque la
compréhension naïve échoue, rien ne garantit que
$\Setabs{X}{\phi(X)}$ existe. De telles définitions semblent donc
\emph{tricher} dangereusement.

Heureusement, elles ne trichent pas ; ou plutôt, si elles
\emph{trichent} telles quelles, un travail honnête permet de les rendre
légitimes. Nous avons annoncé ce travail à la
\olref[sth][z][sep]{prop:intersectionsexist}, en expliquant pourquoi
$\bigcap A$ existe lorsque $A \neq \emptyset$. Détaillons-le
maintenant.

Étant donné l'extensionnalité, si nous essayons de définir $C$ par
$\bigcap\Setabs{X}{\phi(X)}$, nous demandons simplement un objet $C$
qui satisfait :
\begin{equation}\ollabel{bicondelimarbintersection}
	\forall x(x \in C \liff \forall X(\phi(X) \lif x \in X))\tag{*}
\end{equation}
Supposons maintenant qu'il existe \emph{au moins un} ensemble $S$ tel
que $\phi(S)$. Pour établir l'\olref{bicondelimarbintersection}, nous
pouvons alors définir $C$ par \emph{séparation} :
\[
	C = \Setabs{x \in S}{\forall X(\phi(X) \lif x \in X)}.
\]
Nous vous laissons vérifier, à titre d'exercice, que cette définition
donne bien l'\olref{bicondelimarbintersection}.
%  Fixons $x$. Supposons d'abord que $x \in C$, selon la nouvelle
%  définition ; alors $x \in S$ et $\forall X(\phi(X) \lif x \in X)$ ;
%  comme $\phi(S)$, on obtient la condition $\forall X(\phi(X)
%  \lif x \in X)$. Réciproquement, supposons $\forall X(\phi(X) \lif
%  x \in X)$ ; comme $\phi(S)$, nous avons aussi $x \in S$ ; donc
%  $x \in C$, selon la nouvelle définition.
Cette stratégie générale nous permet d'éviter tout recours apparent
à la compréhension naïve pour définir des intersections. Revenons
au cas particulier qui nous a conduits à cette question, celui de
$\closureofunder{f}{o}$. Au début de la démonstration du
\olref[sfr][infinite][dedekind]{closureproperties}, nous avons noté
que $o \in \ran{f}\cup\{o\}$ et que $\ran{f} \cup \{o\}$ est stable
par $f$. Nous pouvons donc définir l'ensemble voulu ainsi :
\[
	\closureofunder{f}{o} = \Setabs{x \in \ran{f} \cup \{o\}}{(\forall X \ni o)(X \text{ est stable par $f$} \lif x \in X)}.
\]

\end{document}
